\begin{explanationbox}
	\textbf{\textcolor{CExplanation}{一句话直觉}}

	函数 $y=kx+b$（$k,b$ 为常数）的图像是直线，斜率 $k$ 决定其倾斜方向和程度，纵截距 $b$ 决定其与 $y$ 轴的交点。

	\medskip
	\textbf{\textcolor{CExplanation}{核心拆解}}
	\begin{description}[style=nextline, leftmargin=2.8em, labelsep=0.5em, itemsep=0.45em]
		\item[\textbf{要点一（斜率公式）：}] 斜率 $k = \frac{y_2 - y_1}{x_2 - x_1} = \frac{\Delta y}{\Delta x}$，表示自变量每变化一个单位，因变量平均变化 $k$ 个单位。
		\item[\textbf{要点二（纵截距）：}] 纵截距 $b = y(0)$，是直线与 $y$ 轴交点的纵坐标，反映函数在零点处的初始值。
		\item[\textbf{要点三（单调性判据）：}] $k$ 的正负直接决定函数在 $\mathbb{R}$ 上的单调性：$k>0$ 递增，$k<0$ 递减，$k=0$ 不变。
	\end{description}

	\medskip
	\textbf{\textcolor{CExplanation}{几何本质（倾斜程度）}}

	斜率 $k$ 的绝对值越大，直线越陡峭；$k>0$ 时直线从左下到右上，$k<0$ 时从左上到右下；$b$ 是直线在 $y$ 轴上的截距（不是横截距）。

	\medskip
	\textbf{\textcolor{CExplanation}{代数意义（平均变化率）}}

	$k = \frac{y_2 - y_1}{x_2 - x_1}$ 正是函数在区间 $[x_1,x_2]$ 上的平均变化率，而 $b$ 是自变量的初始对应值；当 $k=0$ 时函数为常函数，变化率为零。

	\medskip
	\textbf{\textcolor{CExplanation}{考点价值（常见考法）}}
	\begin{itemize}[leftmargin=*, itemsep=0.5em, topsep=0.4em]
		\item \textbf{考法一（求解析式）：} 已知两点或一点与斜率，利用 $k,b$ 公式求解析式。
		\item \textbf{考法二（判断单调性）：} 由 $k$ 符号直接判断函数增减，无需列表或作图。
		\item \textbf{考法三（图像与平移）：} 改变 $k$ 改变倾斜程度，改变 $b$ 上下平移直线。
	\end{itemize}

	\medskip
	\textbf{\textcolor{CExplanation}{顿悟点}}

	斜率公式与平均变化率定义完全一致，为后续学习导数（瞬时变化率）作铺垫；同时 $k$ 的正负与单调性直接对应，体现了“数形结合”的思想。

	\medskip
	\textbf{\textcolor{CExplanation}{使用场景}}
	\begin{itemize}[leftmargin=*, itemsep=0.5em, topsep=0.4em]
		\item \textbf{场景一（点斜式求方程）：} 已知直线过一点 $(x_0,y_0)$ 且斜率为 $k$，可直接写出 $y-y_0 = k(x-x_0)$。
		\item \textbf{场景二（两点式求方程）：} 已知两点 $(x_1,y_1),(x_2,y_2)$，先求斜率再代点得方程。
		\item \textbf{场景三（线性模型分析）：} 在实际问题中，若变量均匀变化，可用一次函数建模，$k$ 代表变化速率，$b$ 代表初始值。
	\end{itemize}
\end{explanationbox}
